\documentclass[a4paper]{article}

\usepackage{array}
\usepackage{fullpage}
\usepackage{enumerate}
\usepackage{hyperref}
\usepackage[all]{hypcap}
\usepackage{listings}
\usepackage{color}
\usepackage{graphicx}

\definecolor{light-gray}{gray}{0.95}
\definecolor{dark-green}{rgb}{0, 0.5, 0}
\definecolor{dark-gray}{gray}{0.4}
\definecolor{Gray}{gray}{0.95}

\lstset{ %
	language = C,                   % choose the language of the code
	basicstyle = \small\ttfamily,   % the size and fonts that are used
	frame = single,                 % adds a frame around the code
	tabsize = 3,                    % sets default tabsize
	breaklines = true,              % sets automatic line breaking
	numbers = left,                 % where to put the line-numbers
	numberstyle = \footnotesize,    % the style of the line-numbers
	backgroundcolor = \color{Gray}, % the background color of the listing
	showstringspaces=false,
	keywordstyle=\color{blue},
}

\newcolumntype{V}{>{\centering\arraybackslash} m{.75\linewidth} }

\title{Assignment 3: MPI}
\author{Jacco Spoelder (s1348493) \and Xeryus Stokkel (s2332795)}

\begin{document}

\maketitle

\section{Exercise 1}
\begin{enumerate}[(a)]
	\item When cutting the interval into equal chunks we obtain the results in \autoref{tbl:prime}. When using multiple processors we can see that there is a discrepancy in the runtime for different processes. The longest runtime can be four times as high as the lowest runtime. The cause of this is quite obvious, if there are two processes then one will work on the interval $[1, \ldots, 2.5 \cdot 10^7]$ while the other will work on the interval $[2.5 \cdot 10^7, \ldots, 5 \cdot 10^7]$. This means that the first process will need to check roughly $\sqrt{\sum_{n=1}^{2.5 \cdot 10^7}n} \approx 1.7 \cdot 10^7$ numbers, the second process needs to check $\sqrt{\sum_{n=2.5\cdot 10^7}^{5\cdot 10^7}n} \approx 3 \cdot 10^7$ numbers in total in the \texttt{isPrime} function. This is almost a doubling of the amount of possible divisors to check and we can see that the runtime is also almost double from \autoref{tbl:prime}. Once the process with the largest numbers in its interval is done the program will also complete after displaying the total number of primes.
	\item We used the master-worker model as our parallelization strategy. The master process sends out small parts of the interval to different workers which will then work on calculating the number of primes. When all work has been distributed a termination signal will be send to each worker and the worker processes will send their results back to the master process upon receiving the termination signal.
	
	Special care has been taken to ensure the ordering of messages to ensure that workers don't accidentally ignore work because they received the termination signal early. This happened during some of the early test runs so we took the time to fix this without having to fundamentally change how our program works.
	
	Because of the way that communication works the master process doesn't actually do any work on finding the number of primes. Therefore there is no data for 1 process in \autoref{tbl:prime2} since it would never complete. From \autoref{tbl:prime2} we can see that the speed-up is a lot higher than in \autoref{tbl:prime}. It should be noted that there is effectively one process less working on the problem than the table states so the achieved speed-up would be higher if all the processes would've been working on the problem. In \autoref{tbl:prime2} we also see that the difference in highest and lowest runtime is about 100ms so the load is a lot better balanced amongst the processors as none of them finishes early.
\end{enumerate}

\begin{table}[h]
	\centering
	\caption{Runtimes for finding the number of primes by splitting the load into equal sized parts. Runtime for the sequential program is 2:22.72, the speed-up is based on comparing to this time.}
	\label{tbl:prime}
	\begin{tabular}{l|r|r|r|r}
		Processes & Lowest time & Highest time & Total time & Speed-up \\ \hline
		 1 &  & 2:19.60 & 2:19.74 & 1.02 \\
		 2 & 50.61 & 1:28.64 & 1:28.78 & 1.61 \\
		 4 & 18.43 & 47.92 & 48.20 & 2.96 \\
		 8 &  6.74 & 24.71 & 26.16 & 5.46 \\
		12 &  3.74 & 16.67 & 19.36 & 7.37 \\
	\end{tabular}
\end{table}

\begin{table}[h]
	\centering
	\caption{Runtimes for finding the number of primes by interleaving the load. Runtime for the sequential program is 2:22.72, the speed-up is based on comparing to this time.}
	\label{tbl:prime2}
	\begin{tabular}{l|r|r|r|r}
		Processes & Lowest time & Highest time & Total time & Speedup \\ \hline
		 2 &    & 2:15.49 & 2:16.68 & 1.04 \\
	 	 4 & 45.53 & 45.62 & 46.85 & 3.04 \\
	 	 8 & 19.49 & 19.59 & 21.06 & 6.78 \\
		12 & 12.41 & 12.55 & 14.26 & 10.01 \\
	\end{tabular}
\end{table}

\section{Contrast stretching}

\section{Mandelbrot fractal}
\begin{enumerate}[(a)]
	\item Runtimes are displayed per process in \autoref{tbl:mandelbrot}. Each column shows the runtime for each process for a given number of processors. The total runtime for the program is displayed at the bottom. We can see that the columns follow a certain symmetry: the first and last process have the lowest running time while the processes in the middle have the highest running time. The difference in running time is very large when the number of processes grows!
	
	This phenomenon is caused by how we distribute the work. We split the image up in $N$ rows (where $N$ is the number of processes) that are of equal size. From \autoref{fig:mandelbrot} we can see that the complexity at the top and the bottom of the fractal is a lot lower than the complexity in the middle. \autoref{tbl:mandelbrotcalc} shows the image split up into 8 sections, listing the number of iterations required to calculate each section. From this we can see that first and last process have a lot less work to do ($x^2 + y^2 > 10$ is satisfied in fewer iterations). A better solution would be to distribute work in such a way that each process has to deal with the same amount of complexity.
	
	\item When we change the program so that there is a master process that delegates work to slave processes then we obtain results found in \autoref{tbl:mandelbrot2}. Just like with \autoref*{tbl:prime2} there is no time for 1 process because the master doesn't do any work itself. As we can see from \autoref{tbl:mandelbrot2} the times for each process are a lot better balanced. Most processes take an equal amount of time to complete, with the exception of the master process. This is because the time to write the image is also taken included in this time so it will always take longer than the threads that only need to do the calculations.
\end{enumerate}

\begin{table}[h]
	\centering
	\caption{Runtimes for generating an Mandelbrot set image using static domain decomposition. Process times are the time that each process takes to calculate its part of the image, while total time is the time that the program took including memory allocation, MPI communication and writing the fractal to a file.}
	\label{tbl:mandelbrot}
	\begin{tabular}{l|r|r|r|r|r}
		Number of processes & 1 & 2 & 4 & 8 & 12 \\ \hline
		process 0  & 40.76 & 20.36 &  1.14 &  0.06 &  0.03 \\
		process 1  &       & 20.40 & 19.19 &  1.10 &  0.12 \\
		process 2  &       &       & 19.25 &  7.18 &  1.04 \\
		process 3  &       &       &  1.17 & 12.03 &  4.36 \\
		process 4  &       &       &       & 12.06 &  6.20 \\
		process 5  &       &       &       &  7.20 &  8.70 \\
		process 6  &       &       &       &  1.12 &  8.71 \\
		process 7  &       &       &       &  0.06 &  6.18 \\
		process 8  &       &       &       &       &  4.37 \\
		process 9  &       &       &       &       &  1.03 \\
		process 10 &       &       &       &       &  0.11 \\
		process 11 &       &       &       &       &  0.04 \\ \hline
		Total      & 41.63 & 21.03 & 20.92 & 13.95 & 10.19
	\end{tabular}
\end{table}

\begin{figure}
	\centering
	\includegraphics[width=.75\textwidth]{mandelbrot/mandelbrot.png}
	\caption{Final render of the Mandelbrot fractal}
	\label{fig:mandelbrot}
\end{figure}

\begin{table}[h]
	\centering
	\caption{8 sections of the Mandelbrot fractal including and the number of iterations required to compute each section. The intensity of a pixel represents the amount of iterations needed to compute its value. A black pixel means no calculations while a white pixel means that the iteration limit has been reached.}
	\label{tbl:mandelbrotcalc}
	\begin{tabular}{Vr}
		\includegraphics[width=.75\textwidth]{mandelbrot/mandelbrotcalc0.png} & 6 123 840 \\
		\includegraphics[width=.75\textwidth]{mandelbrot/mandelbrotcalc1.png} & 183 359 327 \\
		\includegraphics[width=.75\textwidth]{mandelbrot/mandelbrotcalc2.png} & 1 218 769 786 \\
		\includegraphics[width=.75\textwidth]{mandelbrot/mandelbrotcalc3.png} & 2 042 471 058 \\
		\includegraphics[width=.75\textwidth]{mandelbrot/mandelbrotcalc4.png} & 2 046 873 095 \\
		\includegraphics[width=.75\textwidth]{mandelbrot/mandelbrotcalc5.png} & 1 221 147 602 \\
		\includegraphics[width=.75\textwidth]{mandelbrot/mandelbrotcalc6.png} & 184 863 811 \\
		\includegraphics[width=.75\textwidth]{mandelbrot/mandelbrotcalc7.png} & 6 129 020 \\
	\end{tabular}
\end{table}

\begin{table}[h]
	\centering
	\caption{Runtimes for generating an Mandelbrot set image by using the master-slave model. Process time is the time that a process took for MPI communication and fractal calculation while total time includes setup and writing the obtained image to a file.}
	\label{tbl:mandelbrot2}
	\begin{tabular}{l|r|r|r|r}
		 & 2 & 4 & 8 & 12 \\ \hline
		process 0  & 41.26 & 14.05 & 6.30 & 4.16 \\
		process 1  & 40.83 & 13.62 & 5.84 & 3.73 \\
		process 2  &       & 13.62 & 5.84 & 3.73 \\
		process 3  &       & 13.62 & 5.84 & 3.73 \\
		process 4  &       &       & 5.84 & 3.74 \\
		process 5  &       &       & 5.84 & 3.72 \\
		process 6  &       &       & 5.84 & 3.73 \\
		process 7  &       &       & 5.84 & 3.73 \\
		process 8  &       &       &      & 3.73 \\
		process 9  &       &       &      & 3.73 \\
		process 10 &       &       &      & 3.73 \\
		process 11 &       &       &      & 3.73 \\ \hline
		Total      & 42.44 & 15.30 & 7.71 & 5.92
	\end{tabular}
\end{table}

\end{document}